\begin{proposition}[Two rows permit only elementary two-torsion]
\label{B:two-row-torsion}
Let $M$ be a finite monoid whose minimal ideal is a rectangular band
with two rows or two columns. If $BM$ has finite CW homotopy type, then
the torsion subgroup of $H_2(BM;\mathbb Z)$ is a direct sum of copies
of $\mathbb Z/2$. In particular it has no odd-primary torsion and no
element of order four. If every row action fixes the two rows or is
constant, then $H_2(BM;\mathbb Z)$ is torsion-free.
\end{proposition}
\begin{proof}
Pass to the opposite monoid if necessary, and write the minimal ideal
as $\{0,1\}\times J$. Let $\chi_m\in\{0,1,-1\}$ be the action of $m$
on the row augmentation module, and let $\tau_m:J\to J$ be its right
column action. Fix $j_0\in J$. The natural tensor-product description
in Lemma~\ref{lem:B-band-Tor} presents the abelian group
$A=H_2(BM;\mathbb Z)$ by generators $\ell_j$, with $\ell_{j_0}=0$,
and relations
\begin{equation}\label{B:torsion-affine}
 \ell_{\tau_m(j)}=\chi_m\ell_j+b_m,
 \qquad b_m=\ell_{\tau_m(j_0)}.
\end{equation}
The generator $\ell_j$ is the class of
$(e_j-e_{j_0})\otimes(e_1-e_0)$. Composition gives
\[
 \chi_{mn}=\chi_m\chi_n,\qquad
 b_{mn}=b_n+\chi_n b_m.
\]

We first prove that $b_m=0$ whenever $\chi_m=1$. Let
$X=\{m\in M:\chi_m=1\}$. This is a finite submonoid, and
$b:X\to A$ is an additive monoid homomorphism. Its finite image is a
group. The local maximal group of the minimal ideal of $X$ maps onto
this image, by the same finite-semigroup argument as in
Lemma~\ref{A:minimal-group}. Thus it suffices to prove $b_h=0$ for elements
of any local finite subgroup $H\subset X$.

Let $e$ be the identity of $H$. Its row action is the identity, and its
column action is an idempotent projection. Choose a column orbit
$J_0\subseteq\operatorname{im}\tau_e$ for $H$, and let $K\leq H$ be
the stabilizer of a column $j\in J_0$. Inside the corner $eMe$ there is
the finite monoid
\[
 S=H\ \amalg\ (\{0,1\}\times J_0).
\]
The group fixes the two rows and acts transitively on the columns.
The cofiber description in Theorem~\ref{E:group-band-cofiber}
therefore gives
\begin{equation}\label{B:torsion-group-orbit}
 BS\simeq\Sigma\operatorname{hocofib}(BK\longrightarrow BH).
\end{equation}
For completeness, the Borel construction of the join of the trivial
two-point row set and the orbit $H/K$ is
$BH\cup^h_{BK}BH$. Its map to $BH$ is the fold map. The homotopy
cofiber of that fold is the suspension of the homotopy cofiber of
$BK\to BH$, which proves the displayed formula.

Every based map from the right side of
\eqref{B:torsion-group-orbit} to a finite connected CW complex is
null-homotopic. Indeed Miller's full mapping-space theorem says that
the based mapping spaces from $BH$ and $BK$ to such a target are
weakly contractible. Applying based mapping spaces to the cofiber
sequence and then looping gives the assertion. Apply it to a finite
CW model of $BM$. The corner inclusion $BS\to BM$ is consequently
null-homotopic.

The degree-two band-Tor formula is natural for the inclusion of $S$
into the corner and then into $M$. It identifies
\[
 H_2(BS;\mathbb Z)
   \cong \bigl(\ker(\mathbb ZJ_0\to\mathbb Z)\bigr)_H.
\]
To interpret this naturality using only unital monoid maps, adjoin a
fresh identity to $S$. The new monoid $S^+$ maps unitally to $M$,
and its old-identity corner is $S$ with the usual one-sided corner
equivalence $BS\simeq BS^+$. Its minimal band and balanced tensor
presentation are unchanged. Thus the ordinary unital naturality of
Lemma~\ref{lem:B-band-Tor} gives exactly the stated corner map.
The coinvariant class of $e_{\tau_h(j)}-e_j$ maps to
$\ell_{\tau_h(j)}-\ell_j=b_h$ in $A$, by
\eqref{B:torsion-affine}. Since the inclusion is null-homotopic,
$b_h=0$. Surjectivity from the local group proves $b_m=0$ for every
$m\in X$.

If there are no elements with $\chi_m=-1$, the remaining relations
in \eqref{B:torsion-affine} only identify pairs of column values:
positive-character maps identify $\ell_{\tau_m(j)}$ with $\ell_j$,
and zero-character maps identify all values in their image. Together
with $\ell_{j_0}=0$, this presents a free abelian group.

Otherwise all negative-character constants have the same value $c$.
Indeed for $\chi_m=\chi_n=-1$ the product has positive character, so
$0=b_{mn}=b_n-b_m$. The presentation is now a signed graph on the
column variables: positive edges impose $\ell_v=\ell_w$, negative
edges impose $\ell_v+\ell_w=c$, and $\ell_{j_0}=0$ anchors one vertex.
An extra variable $c$ with its defining relation can be used without
changing the group. Each negative action supplies the negative edge
from $j_0$ to $\tau_m(j_0)$, so this defining relation is already one
of the signed-graph relations.

In a connected component with no cycle containing an odd number of
negative edges, all its variables are either $x$ or $c-x$, for one
free variable $x$. If there is an odd signed cycle, the additional
relation is exactly $2x=c$. The anchored component has values $0,c$;
an odd signed cycle there forces $c=0$. These conclusions follow by
choosing a spanning tree in each component and then examining each
remaining edge, so they describe the complete presentation.
If $c=0$, the only torsion relations are $2x=0$. If $c\ne0$ and there
are odd components, eliminate $c$ using one relation $2x_1=c$; every
other such relation becomes $2(x_i-x_1)=0$. With no odd components
the presentation is free. Thus the only possible torsion summands
are $\mathbb Z/2$, as claimed.
\end{proof}
